\section{Model Implementation}
\label{sec:model}

\subsection{Architecture}

The generator is a decoder-only transformer~\cite{vaswani2017attention}. The structural
backbone --- causal self-attention, weight tying between the input embedding and the language
model head, the $\mathcal{N}(0, 0.02)$ initialisation scheme with the residual-scaled
projection variance \cite{radford2019gpt2} and the \texttt{c\_attn}/\texttt{c\_proj} naming
convention --- follows nanoGPT~\cite{karpathy2022nanogpt}, a clean reimplementation of the
GPT-2 design. Three departures from that baseline were introduced to align the architecture
with components that are now standard in modern decoder-only models such as
LLaMA~\cite{touvron2023llama}: rotary positional embeddings, RMSNorm, and a SwiGLU
feed-forward block. Two further additions, stochastic depth and a key/value cache, address
training stability and inference latency respectively. Each is summarised below and laid out
schematically in Figure~\ref{fig:architecture}.

\begin{figure}[h]
\centering
\begin{tikzpicture}[node distance=0.4cm]

    % Input
    \node[data] (tokens) {Token IDs \\ \footnotesize $(B, T)$};
    \node[block, above=of tokens] (embed) {Token embedding \\ \footnotesize tied to LM head};

    % Block (single)
    \node[draw=black!50, dashed, rounded corners=4pt, fill=gray!3,
          minimum width=8.5cm, minimum height=5.6cm,
          above=of embed, yshift=0.3cm] (blockbox) {};
    \node[smallnote, anchor=south east] at ($(blockbox.south east)+(-0.2,0.15)$) {$\times N$ blocks};

    \node[block, above=1.4cm of embed, fill=blue!4] (norm1) {RMSNorm};
    \node[block, above=of norm1, fill=blue!8] (attn) {Causal multi-head attention \\ \footnotesize with RoPE on Q,K (+ KV cache)};
    \node[block, above=of attn, fill=blue!4] (norm2) {RMSNorm};
    \node[block, above=of norm2, fill=blue!8] (ffn) {SwiGLU FFN \\ \footnotesize hidden $\approx \tfrac{8}{3}d$};

    % Stochastic depth note
    \node[smallnote, right=0.5cm of attn] (sd1) {+ residual\\(stoch.\ depth)};
    \node[smallnote, right=0.5cm of ffn]  (sd2) {+ residual\\(stoch.\ depth)};

    % Final norm + head
    \node[block, above=1.6cm of ffn] (normf) {RMSNorm};
    \node[block, above=of normf] (head) {LM head \\ \footnotesize logits $(B, T, V)$};

    % Edges
    \draw[arrow] (tokens) -- (embed);
    \draw[arrow] (embed) -- (norm1);
    \draw[arrow] (norm1) -- (attn);
    \draw[arrow] (attn) -- (norm2);
    \draw[arrow] (norm2) -- (ffn);
    \draw[arrow] (ffn) -- (normf);
    \draw[arrow] (normf) -- (head);

\end{tikzpicture}
\caption{Forward pass of the custom transformer. Pre-norm RMSNorm wraps causal multi-head
attention (with RoPE applied to queries and keys) and a SwiGLU feed-forward block; the two
sub-blocks are skipped with probability $p$ during training (stochastic depth). $V = 50\,261$
tokens, $d = 512$, $N = 6$.}
\label{fig:architecture}
\end{figure}

\paragraph{Rotary positional embeddings.}
Position information is injected by rotating the query and key projections by an angle that
depends on the absolute position of each token~\cite{su2024roformer}. Unlike learned absolute
embeddings, RoPE makes the relative position between two tokens an explicit feature of the
attention dot product, which extrapolates better to sequence lengths not seen during training
and is the dominant choice in current open-weight models.

\paragraph{RMSNorm.}
Each pre-norm step uses root mean square layer normalisation~\cite{zhang2019rmsnorm} rather
than full LayerNorm. RMSNorm drops the mean-centring term, which has been observed to have
limited effect on quality while removing one reduction operation per step.

\paragraph{SwiGLU.}
The feed-forward block uses the SwiGLU gating variant~\cite{shazeer2020glu}, $\mathrm{FFN}(x)
= (x W_1 \odot \mathrm{SiLU}(x W_2)) W_3$. To keep the parameter count comparable to a GELU
$4d$ block, the inner dimension is set to approximately $\tfrac{8}{3}d$. For $d=512$ this
gives $\mathrm{round}(8/3 \cdot 512) = 1\,365$, where the GELU variant of the same model
would use $4d = 2\,048$.

\paragraph{Stochastic depth.}
During training, each residual sub-block (attention or FFN) is independently dropped with a
small linearly increasing probability across depth~\cite{huang2016stochastic}. This regularises
deeper blocks more aggressively than shallow ones and was added because the model has
relatively few examples per parameter once the synthetic corpus is restricted to in-domain
chunks.

\paragraph{KV cache.}
At inference time, autoregressive decoding recomputes the same key and value projections for
every previously generated token. The implementation caches them so each new step only
projects the most recent token, reducing the per-step cost from $O(T^2)$ to $O(T)$ when
sampling long outputs (notably for the \emph{review} task).

\paragraph{Hyperparameters.}
Table~\ref{tab:hyperparams} summarises the architectural settings. The model has roughly
26\,M parameters with weight tying.

\begin{table}[h]
\centering
\begin{tabular}{ll}
\toprule
\textbf{Parameter} & \textbf{Value} \\
\midrule
Layers ($N$)                & 6 \\
Attention heads             & 8 \\
Embedding dimension ($d$)   & 512 \\
FFN hidden dim (SwiGLU)     & 1\,365 ($\approx \tfrac{8}{3}d$) \\
FFN hidden dim (GELU var.)  & 2\,048 ($4d$) \\
Context length              & 512 tokens \\
Positional encoding         & RoPE \\
Normalisation               & RMSNorm \\
Vocabulary                  & GPT-2 BPE + 4 task tokens (50\,261) \\
\bottomrule
\end{tabular}
\caption{Architecture settings of the custom transformer. The FFN dimension is
variant-dependent.}
\label{tab:hyperparams}
\end{table}

Attention is computed via \texttt{torch.nn.functional.scaled\_dot\_product\_attention}, which
dispatches to a FlashAttention kernel when one is available and falls back to a memory-safe
reference implementation otherwise.

\subsection{Tokenisation and task tokens}
\label{sec:tokenizer}

The four task formats --- explanation, quiz, review, flashcard --- are signalled to the model
by four dedicated special tokens that prefix every example. Reusing GPT-2's existing BPE
vocabulary avoided retraining a tokeniser from scratch, but four IDs had to be appended to
the vocabulary:

\begin{center}\small
\begin{tabular}{lll}
\toprule
\texttt{<|explain|>} & ID 50\,257 & free-form explanation \\
\texttt{<|quiz|>}    & ID 50\,258 & multiple-choice quiz \\
\texttt{<|review|>}  & ID 50\,259 & long-answer + scored review \\
\texttt{<|flashcard|>} & ID 50\,260 & front/back card \\
\bottomrule
\end{tabular}
\end{center}

The implementation rebuilds the tiktoken \texttt{gpt2} encoding with the four extra entries
in the \texttt{special\_tokens} table, preserving the original 50\,256 BPE merges and
the \texttt{<|endoftext|>} marker (ID 50\,256) used as end-of-sequence. The model's embedding
table and language model head are sized to 50\,261 to match. Conditioning on a task token in
this way is a lightweight alternative to training separate heads per task and lets a single
forward pass cover all four activities; the same approach is used in recent
instruction-tuned models such as T5~\cite{raffel2020t5}.

\subsection{Synthetic training data}
\label{sec:training_data}

A model trained from random initialisation on the public QA corpora alone would learn to
answer general machine-learning questions, but it would not learn to emit the four task
formats above. To produce supervision for those formats, an instruction-tuned local model,
\texttt{Qwen2.5-7B-Instruct}~\cite{yang2024qwen25}, was used to generate task-specific
training examples from the chunks indexed in Section~\ref{sec:data}. Running the data
generator locally avoided API costs and let the same chunk be expanded into all four task
types in a single pass over the corpus. Figure~\ref{fig:datagen} sketches the pipeline.

\begin{figure}[h]
\centering
\begin{tikzpicture}[node distance=0.5cm and 0.7cm]
    \node[data] (chunks) {Indexed chunks \\ \footnotesize \texttt{chunks.json}};
    \node[block, right=of chunks] (qwen) {Qwen2.5-7B \\ \footnotesize 4 task prompts};

    \node[data, right=1.3cm of qwen, yshift=1.5cm]   (e) {\texttt{explanation}};
    \node[data, right=1.3cm of qwen, yshift=0.5cm]   (q) {\texttt{quiz}};
    \node[data, right=1.3cm of qwen, yshift=-0.5cm]  (r) {\texttt{review}};
    \node[data, right=1.3cm of qwen, yshift=-1.5cm]  (f) {\texttt{flashcard}};

    \node[data, right=1.5cm of q, yshift=-0.5cm] (jsonl) {\texttt{train.jsonl}};

    \draw[arrow] (chunks) -- (qwen);
    \draw[arrow] (qwen.east) -- (e.west);
    \draw[arrow] (qwen.east) -- (q.west);
    \draw[arrow] (qwen.east) -- (r.west);
    \draw[arrow] (qwen.east) -- (f.west);
    \draw[arrow] (e.east) -- (jsonl.north west);
    \draw[arrow] (q.east) -- (jsonl.west);
    \draw[arrow] (r.east) -- (jsonl.west);
    \draw[arrow] (f.east) -- (jsonl.south west);
\end{tikzpicture}
\caption{Synthetic data generation. Each indexed chunk yields one example per task type.
Resume state is tracked by SHA-1 hash of the chunk text, so an interrupted run continues at
chunk granularity rather than chapter granularity.}
\label{fig:datagen}
\end{figure}

A single line of \texttt{train.jsonl} stores the rendered training string under
\texttt{"text"} together with the structured fields used for evaluation. The example below is
a \emph{quiz} record (truncated for readability):

\begin{quote}\footnotesize\ttfamily
\{ \\
\hspace*{1em}"type": "quiz", \\
\hspace*{1em}"text": "<|quiz|>$\backslash$nContext: At this point, if you want to further \\
\hspace*{2em}improve performance$...\backslash$nQuiz: \{$\backslash$"question$\backslash$": $...$, \\
\hspace*{2em}$\backslash$"options$\backslash$": [$...$], $\backslash$"correct\_index$\backslash$": 2, \\
\hspace*{2em}$\backslash$"explanation$\backslash$": $...$\}<|endoftext|>", \\
\hspace*{1em}"quiz": \{ "question": ..., "options": [...], "correct\_index": 2 \}, \\
\hspace*{1em}"source": ".../ch09\_convnet-architecture-patterns.md", \\
\hspace*{1em}"chunk\_hash": "1772b1db..." \\
\}
\end{quote}

The structured copy of each field is preserved alongside the rendered string so that
format-validity metrics (Section~\ref{sec:evaluation}) can be computed without re-parsing the
generated string. The \texttt{chunk\_hash} provides a stable resume key: an earlier
implementation that tracked progress by chapter source path skipped every chunk in a chapter
once the first one was processed; switching to a SHA-1 hash of the chunk text fixed this and
made interrupted runs safely restartable.

For the explanation and flashcard tasks, Qwen is asked to write a request and the
corresponding response in the voice of a teaching assistant. For \emph{quiz}, the prompt
constrains the output to a strict JSON schema with four options and a \texttt{correct\_index}
field. For \emph{review}, Qwen first generates a question and a plausible (and deliberately
imperfect) student answer, then writes a graded review of that answer with a numeric score
and a reference answer. The reference answer is what the model is later judged against.

The validity of the generated data is checked by a \texttt{--validate} pass that re-parses
each line and rejects entries with malformed task fields, missing context, or a length below
a minimum threshold; rejected entries are not used in training.

\subsection{Training procedure}

The model is trained with the standard autoregressive cross-entropy objective over the
rendered \texttt{text} field. Loss is computed on every token (no prefix masking), since the
task prefix and context are themselves part of the in-distribution sequences the model needs
to learn to recognise. Label smoothing of $0.1$ is applied to discourage extreme logit
values and improve calibration on the JSON-formatted tasks where overconfident wrong tokens
break parsing.

\paragraph{Optimisation.}
The optimiser is AdamW~\cite{loshchilov2019decoupled} with $\beta_1 = 0.9$, $\beta_2 = 0.95$,
weight decay $0.1$, and a peak learning rate of $3 \times 10^{-4}$ on a cosine schedule with
200 linear warmup steps. Parameter groups are split by parameter \emph{name}, not by tensor
shape: every embedding-like parameter is excluded from weight decay regardless of its
dimensionality. This is a small but important correction relative to the common shape-based
heuristic, which silently regularises the embedding table on tied-weight models. Gradient
norms are clipped to $1.0$.

\paragraph{Hardware utilisation.}
Training runs on a single NVIDIA RTX 3090. Forward and backward passes use \texttt{bfloat16}
autocasting; the master weights stay in \texttt{float32}. The model is wrapped in
\texttt{torch.compile} to fuse the attention and FFN kernels. With a per-step batch size of
16 and gradient accumulation of 4, the effective batch size is 64 sequences of 512 tokens.

\paragraph{Stabilisation tricks.}
An exponential moving average (EMA) of the parameters is maintained alongside the live
weights, with a decay of $0.999$~\cite{tarvainen2017mean}. The EMA copy is used for
validation and final inference; this produces visibly cleaner training curves and reduces
the variance of the validation loss across seeds. Early stopping on validation loss with a
patience of 3 epochs prevents overfitting on the smaller in-domain split.

\paragraph{Logging.}
Training and validation curves are streamed to Weights \& Biases. In addition to standard
loss and perplexity, the validation loop computes a format-compliance percentage: 16 examples
per task type are sampled with EMA weights and decoded, and the proportion that parse to a
valid JSON object (for \emph{quiz}) or contain the expected delimiters (for \emph{review}
and \emph{flashcard}) is logged. This second axis turned out to be considerably more
informative than perplexity for selecting checkpoints, since perplexity decreased smoothly
even when the JSON output was syntactically broken.
